% Limitations (Task 42): five grouped paragraphs (plan 32); handoff/37_narrative_v2.md section 3.5 applied; items
% closed by Stage B removed.
\section{Limitations}
\label{sec:limitations}

\paragraph{Threat model scope.}
All attacks changed flow features, not packets, and a feature vector that our checker accepts may still not correspond
to traffic that an attacker can produce.
The attackers knew the model weights and, for A6, A6x, A8, A8x and the attackers of the controls, the detector; this is
a strong assumption for evasion and a reasonable one for evaluating detectors, but results for weaker attackers may
differ.
The split into \nControls{} controls, \nDerived{} derived and \nImmutable{} immutable features, the budget rule and the
class-family freezes were designed by hand from feature definitions, and other reasonable choices could change attack
success and the attention measurements.
At $\epsilon = \paramEpsHead$ a control may change by up to its full clean value, and the random control A0 already
flipped \azeroBenignShareVzeroSeventeen{} (2017, V0) and \azeroBenignShareVzeroEighteen{} (2018) of the row observations
to benign, so the results at $\epsilon = 0.5$ and $\epsilon = \paramEpsHead$ should be read together
(Appendix~\ref{sec:supp}).
We did not attack IG at an equal budget, so the stability of attention is not compared with that of an attribution
method.

\paragraph{Attacker strength.}
A8x used surrogate gradients (soft ranks, temperature annealed from \axTauStart{} to \axTauEnd) for the
\nRankScores{} rank-based consistency scores, which have no informative exact gradient; a better surrogate or a
gradient-free search on these terms could find further undetected rows, which could only raise the share for D$_+$.
The equal-effort protocol equalises candidates, steps and settings, not optimiser quality, and we did not measure
whether A6 is as strong an attack on D as A8x is on D$_+$; the A8x rows evaded D about as often as the A6 rows aimed at
it (Section~\ref{sec:a6x}), but this does not settle the question.
The attackers of a1 and a2 used stand-ins (median $R^2$ \rtwoMedAoneSeventeen{} and \rtwoMedAtwoSeventeen{} on 2017,
\rtwoMedAoneEighteen{} and \rtwoMedAtwoEighteen{} on 2018), like A8, so their results may understate how easily these
controls can be evaded, and their comparison with D$_+$ under A8x is not at equal attacker strength.
The attackers of all controls kept the selection rule of A6; A6x showed that this rule and the rule of A8x give the
same flipped and undetected status for D, which was not tested for the controls.
The attribution-only control b was not attacked adaptively: its run exceeded the memory of the graphics card and was
stopped before any row was written, and an exact-objective attacker for b was not designed, so its family is untested
against an adaptive attacker.

\paragraph{Detector scope.}
D is a linear model on four scores and D$_+$ adds \nConsistency{} linear terms; non-linear detectors might use the same
scores differently.
Control a1 uses the true class of every row, also at test time, and keeps the class-conditional distribution of
attention; its realised test FPR on 2017 (\ablFprAoneSeventeen) was slightly higher than that of D$_+$
(\ablFprDplusSeventeen; V1, $\epsilon = \paramEpsHead$).
The validation AUROC of every detector is in-sample.
The consistency scores of D$_+$ include IG, which costs much more than attention (on 2017, \msPerRowIGRange{} ms per
row for IG with \apIGSteps{} steps against \msPerRowAttnRange{} ms for attention); we did not measure the end-to-end
latency of the detectors.

\paragraph{Data.}
The 2018 data have no timestamps, so we split at random.
Test rows within \paramNDThr{} of a training row were common in both datasets (Section~\ref{sec:rq3}); the filter was
applied after the experiments, and models and detectors were trained with these rows present.
Known design and error issues exist in benchmark NIDS datasets more generally \cite{flood2024,lanvin2022}.
Web attack had \nWebTestEighteen{} test rows and \nWebValEighteen{} validation explanation rows on 2018, so its per-class
estimates are uncertain; FTP-Patator and DoS Slowhttptest were removed from 2018 and PortScan does not occur in it, so
all comparisons use the \nCommonClasses{} common classes.
On 2018 D alone left only \axHeadVoneEighteenD{} of the source rows successful and undetected, so the detector
comparisons there have little room to differ.
The 2018 models ran to the maximum number of epochs without early stopping, so they may not have converged
(Appendix~\ref{sec:appendix}).

\paragraph{Statistics and scope.}
We used \nSeeds{} seeds per dataset, and several results varied in sign across seeds; in the equal-effort comparison
against A8x, \axHeadOppSeedsSeventeen{} seed on 2017 and \axHeadOppSeedsEighteen{} on 2018 had the opposite sign.
Our bootstrap intervals resample source rows within seed (equal-effort union) or pooled over seeds (paired TPR
differences), reuse the same source rows in every seed and do not capture seed variance, so they do not show that an
effect holds across seeds.
The attackers A8x and A6x were designed after the earlier results were known; their outcome rules and both amendments
were recorded before any of their attack rows existed (Section~\ref{sec:protocol}).
We studied two datasets and one model family, without human evaluation or deployment; faithfulness was measured by
deletion and insertion curves, not by human judgement, and we ran a parameter-randomisation check
(Appendix~\ref{sec:rand}) but no label-randomisation check.
